\section{Overview}

The timeline presented here outlines the anticipated trajectory of the dissertation research and writing from anticipated proposal approval (December 2025) through final defense (December 2026). The schedule is designed to accommodate the iterative, experimental nature of practice-based research-creation while maintaining clear milestones for progress and accountability.

As a research-creation dissertation, this timeline embraces the recursive relationship between making, documenting, and writing that characterizes practice-based inquiry. Rather than separating "research," "creation," and "writing" into discrete phases, this approach recognizes that theoretical insights emerge through material practice, and documentation occurs concurrently with making. \textit{Substantial portions of the dissertation have already been developed through the proposal process, establishing a foundation for continued writing-as-making throughout the research period.}

The timeline is structured around three overlapping phases that mirror the cyclical nature of Speculocultural Technopoiesis while allowing for the experimental pivots and productive failures inherent to creative technical work. Built-in flexibility accounts for the unpredictable nature of collaborative practice, technical experimentation, and the emergence of unanticipated insights that may redirect the work.

\textbf{Target Defense Date}: December 2026 (Winter Quarter, AY 2026-2027)

\textbf{Total Research Period}: Approximately 13 months (November 2025 – December 2026)

\section{Timeline Visualization}

\subsection*{Project Timeline at a Glance}
% GNATT CHART
\begin{ganttchart}[
	hgrid,
	vgrid,
	x unit=0.8cm,
	y unit title=0.6cm,
	y unit chart=0.5cm,
	title height=1
	]{1}{14}
	\gantttitle{2025-2026}{14} \\
	\gantttitle{Nov}{1}
	\gantttitle{Dec}{1}
	\gantttitle{Jan}{1}
	\gantttitle{Feb}{1}
	\gantttitle{Mar}{1}
	\gantttitle{Apr}{1}
	\gantttitle{May}{1}
	\gantttitle{Jun}{1}
	\gantttitle{Jul}{1}
	\gantttitle{Aug}{1}
	\gantttitle{Sep}{1}
	\gantttitle{Oct}{1}
	\gantttitle{Nov}{1}
	\gantttitle{Dec}{1} \\
	
	\ganttbar{Phase 1: Foundation}{1}{5} \\
	\ganttbar{Phase 2: Deep Making}{4}{9} \\
	\ganttbar{Phase 3: Synthesis}{9}{14} \\
	
	\ganttbar[bar/.append style={fill=purple!30}]{URET Cycle 1}{2}{3} \\
	\ganttbar[bar/.append style={fill=purple!30}]{URET Cycle 2}{5}{6} \\
	\ganttbar[bar/.append style={fill=purple!30}]{URET Cycle 3}{7}{8} \\
	\ganttbar[bar/.append style={fill=purple!30}]{URET Cycle 4}{10}{11} \\
	
	\ganttbar[bar/.append style={fill=gray!20}]{Documentation (ongoing)}{1}{14} \\
	
	\ganttmilestone{Proposal Defense}{2} \\
	\ganttmilestone{Committee Check-in 1}{5} \\
	\ganttmilestone{Committee Check-in 2}{9} \\
	\ganttmilestone{Defense}{14}
\end{ganttchart}

\section{Milestones}

\subsection*{Phase 1: Foundation and Initial Iterations (November 2025 – March 2026)}
\textit{4-5 months | Focus: Technical infrastructure, initial making cycles, establishing collaborative relationships}

\textbf{November – December 2025}
\begin{itemize}[noitemsep]
	\item Proposal defense and approval
	\item Establish collaborative relationships and partnerships (identify 2-3 potential co-creators)
	\item Set up technical infrastructure (hardware, software environments, documentation systems)
	\item Begin first making cycle: initial experiments with cultural sonic encoding
	\item Concurrent documentation: field notes, process videos, technical logs
	\item \textit{Deliverable}: Technical environment established; proposal completed and presented
\end{itemize}

\textbf{January – February 2026}
\begin{itemize}[noitemsep]
	\item First URET cycle
	\item Iterative development of sonic/computational prototypes
	\item Weekly documentation and reflective writing
	\item Initial drafting: expanding methodology chapter with emergent insights
	\item \textit{Deliverable}: First working prototype(s); draft literature review additions (+15-20 pages)
\end{itemize}

\textbf{March 2026}
\begin{itemize}[noitemsep]
	\item Committee check-in \#1: Share progress, prototypes, preliminary findings
	\item Reflective pause: analyze first iteration, identify productive directions
	\item Begin Methodology chapter expansion (plus Literature Review expansion) based on emergent theoretical needs
	\item \textit{Deliverable}: Committee presentation; revised timeline/research direction if needed
\end{itemize}

\subsection*{Phase 2: Deep Making and Collaborative Practice (April – August 2026)}
\textit{5 months | Focus: Multiple making cycles, collaborative sessions, sustained documentation, chapter drafting}

\textbf{April – May 2026}
\begin{itemize}[noitemsep]
	\item Second and third URET cycles: building on first iteration
	\item Collaborative sonic computation sessions with co-creators
	\item Expanded technical experiments: hardware/software integration, sculptural elements
	\item Concurrent writing: drafting Chapter 4 (Findings/Practice Documentation)
	\item \textit{Deliverable}: 2-3 prototypes/artifacts at various stages; draft chapter (25-30 pages)
\end{itemize}

\textbf{June – July 2026}
\begin{itemize}[noitemsep]
	\item Fourth URET cycle: refinement and transmitting phase emphasis
	\item Documentation curation: select representative examples, prepare visual materials
	\item Test transmitting practices: share work with broader communities, gather feedback
	\item Continue drafting: Chapter 5 (Analysis/Reflection on Practice)
	\item \textit{Deliverable}: Documented collaborative sessions; draft chapter (20-25 pages)
\end{itemize}

\textbf{August 2026}
\begin{itemize}[noitemsep]
	\item Committee check-in \#2: Share substantial work and drafted chapters
	\item Assess progress toward completion; identify remaining gaps
	\item Begin consolidating documentation into dissertation format
	\item \textit{Deliverable}: Committee review of 2-3 chapters; feedback integration plan
\end{itemize}

\subsection*{Phase 3: Synthesis, Writing, and Defense Preparation (September – December 2026)}
\textit{4 months | Focus: Final making cycle, synthesis writing, defense preparation}

\textbf{September – October 2026}
\begin{itemize}[noitemsep]
	\item Final URET cycle: synthesis and culminating work
	\item Intensive writing: complete remaining chapters (Introduction, Conclusion, Implications)
	\item Integration: weave practice documentation throughout dissertation
	\item Finalize visual documentation: figures, images, diagrams, supplementary materials
	\item \textit{Deliverable}: Complete draft of dissertation (all chapters)
\end{itemize}

\textbf{November 2026}
\begin{itemize}[noitemsep]
	\item Submit complete draft to committee for review
	\item Prepare defense presentation and supporting materials
	\item Create defense artifacts: demonstration, portfolio, or installation documentation
	\item Revisions based on committee feedback
	\item \textit{Deliverable}: Revised final dissertation; defense presentation materials
\end{itemize}

\textbf{December 2026}
\begin{itemize}[noitemsep]
	\item Final committee review and approval to defend
	\item Dissertation defense (target: December 2026)
	\item Final revisions and formatting
	\item Submission to Graduate School
	\item \textit{Deliverable}: Successful defense; approved dissertation
\end{itemize}

\subsection*{Concurrent Activities (Throughout Research Period)}

\textbf{Ongoing Documentation Practices}
\begin{itemize}[noitemsep]
	\item Weekly reflective writing and field notes
	\item Process videos and photographic documentation
	\item Technical logs and code repositories (GitHub)
	\item Collaborative session recordings and transcriptions
\end{itemize}

\textbf{Writing-as-Making}
\begin{itemize}[noitemsep]
	\item Continuous writing integrated with practice (not saved for end)
	\item Regular expansion of theoretical framework as concepts emerge
	\item Iterative chapter development in response to making
\end{itemize}

\textbf{Committee Engagement}
\begin{itemize}[noitemsep]
	\item Informal updates and consultations as needed
	\item Two formal committee check-ins (March, August)
	\item Draft chapter submissions for feedback (ongoing)
\end{itemize}

\subsection*{Contingency Planning}

This timeline builds in flexibility for the following contingencies:

\begin{itemize}[noitemsep]
	\item \textbf{Technical failures}: Hardware/software issues are expected in experimental work; alternative approaches and backup systems will be maintained
	\item \textbf{Collaborative scheduling}: If co-creators become unavailable, the autoethnographic solo practice can continue independently
	\item \textbf{Productive pivots}: If early iterations suggest more promising directions, the timeline allows for methodological adjustments (particularly in Phase 1-2)
	\item \textbf{Writing pace}: Extensive proposal development has established writing momentum; concurrent documentation ensures material is ready for integration
\end{itemize}

\subsection*{Alignment with Institutional Requirements}

\begin{itemize}[noitemsep]
	\item \textbf{Committee formation}: Complete (proposal stage)
	\item \textbf{Proposal defense}: December 2025 (Phase 1 start)
	\item \textbf{Candidacy status}: Already achieved upon Candidacy pass/approval
	\item \textbf{Committee check-ins}: Two formal meetings (March, August 2026)
	\item \textbf{Complete draft submission}: November 2026 (4 weeks before defense)
	\item \textbf{Final defense}: December 2026
	\item \textbf{Graduation}: Fall 2026 or Winter 2027 (semester system may be in place by this point)
\end{itemize}

I understand this timeline is ambitious but I also believe it is achievable given the substantial groundwork laid in the proposal, the ongoing nature of the technical practice, and the research-creation approach that integrates making, documenting, and writing as simultaneous rather than sequential activities. The cyclical structure honors the Speculocultural Technopoietic framework while maintaining clear forward momentum toward completion.

\subsection*{Timeline Summary Table}

\begin{table}[H]
	\centering
	\small
	\begin{tabular}{|l|p{3cm}|p{4cm}|p{5cm}|}
		\hline
		\textbf{Phase} & \textbf{Timeframe} & \textbf{Key Activities} & \textbf{Deliverables} \\
		\hline
		\rowcolor{blue!10}
		\textbf{Phase 1} & Nov 2025 – Mar 2026 (5 months) & 
		\begin{itemize}[nosep,leftmargin=*]
			\item Proposal defense
			\item Technical setup
			\item First URET cycle
			\item Initial collaborations
		\end{itemize} & 
		\begin{itemize}[nosep,leftmargin=*]
			\item Working prototypes
			\item Methodology drafts
			\item Committee check-in 1
		\end{itemize} \\
		\hline
		\rowcolor{green!10}
		\textbf{Phase 2} & Apr – Aug 2026 (5 months) & 
		\begin{itemize}[nosep,leftmargin=*]
			\item Multiple URET cycles
			\item Collaborative sessions
			\item Deep making
			\item Chapters 4-5 drafting
		\end{itemize} & 
		\begin{itemize}[nosep,leftmargin=*]
			\item 2-3 prototypes
			\item Draft chapters
			\item Committee check-in 2
		\end{itemize} \\
		\hline
		\rowcolor{orange!10}
		\textbf{Phase 3} & Sep – Dec 2026 (4 months) & 
		\begin{itemize}[nosep,leftmargin=*]
			\item Final URET cycle
			\item Synthesis writing
			\item Defense prep
			\item Revisions
		\end{itemize} & 
		\begin{itemize}[nosep,leftmargin=*]
			\item Complete draft
			\item Defense materials
			\item Final dissertation
		\end{itemize} \\
		\hline
		\rowcolor{gray!10}
		\textbf{Ongoing} & Throughout & 
		\begin{itemize}[nosep,leftmargin=*]
			\item Documentation
			\item Writing-as-making
			\item Committee engagement
		\end{itemize} & 
		Continuous outputs \\
		\hline
	\end{tabular}
	\caption{Phase-based timeline with overlapping activities and iterative cycles.}
	\label{tab:timeline}
\end{table}
